\begin{frame}{Перенос и основной показатель оценки}
\textbf{Перенос:} начальные параметры временного модуля и их адаптация по обучающим данным человека.
\SlideGap
\textbf{Основной показатель:} корень из средней квадратичной ошибки (RMSE) амплитуды кривой роста в микровольтах.
\[
 E_i=\sqrt{\frac{1}{m_i}\sum_{j=1}^{m_i}(\widehat A_{ij}-A_{ij})^2},
 \qquad E=\frac{1}{N}\sum_{i=1}^{N}E_i.
\]
\noindent $E_i$ --- ошибка для участника $i$; $E$ --- средняя ошибка; $N$ --- число участников; $m_i$ --- число точек кривой роста участника $i$; $A_{ij}$ и $\widehat A_{ij}$ --- измеренная и расчётная амплитуды в микровольтах.
\SlideGap
\begin{itemize}
\item Ошибки участников усредняются с равным весом. Повторные измерения не увеличивают число независимых участников.
\item Настройка параметров, масштаба времени и калибровка выполняются внутри обучающих разбиений.
\item Допуски и объём выборки будут определены по пригодным данным до независимой проверки.
\end{itemize}
\end{frame}
